\documentclass[10pt,aspectratio=169]{beamer}
% Preámbulo modular para presentaciones Beamer (Modelación Estadística)
% Diseño y paleta institucional del Tecnológico de Monterrey

\documentclass[aspectratio=169,xcolor={dvipsnames,table}]{beamer}

\usepackage[utf8]{inputenc}
\usepackage[spanish,mexico]{babel}
\usepackage{amsmath,amssymb,amsfonts,mathtools}
\usepackage{graphicx}
\usepackage{listings}
\usepackage{booktabs}
\usepackage{tikz}
\usetikzlibrary{matrix,arrows,decorations.pathmorphing,shapes.geometric,positioning}

% Paleta Institucional Tecnológico de Monterrey
\definecolor{TecAzul}{HTML}{0039A6}       % Azul TEC: color primario institucional
\definecolor{TecAzulOscuro}{HTML}{1A2E51} % Azul marino oscuro
\definecolor{TecRojo}{HTML}{EC2661}       % Rojo de acento (paleta miTec)
\definecolor{TecGris}{HTML}{646464}       % Gris neutro para comentarios y subtítulos
\definecolor{TecFondoBloque}{HTML}{F4F6F9}% Fondo suave para bloques

% Configuración de Tema Beamer
\usetheme{Boadilla}
\usecolortheme[named=TecAzulOscuro]{structure}
\setbeamercolor{alerted text}{fg=TecRojo}
\setbeamercolor{palette primary}{bg=TecAzulOscuro,fg=white}
\setbeamercolor{palette secondary}{bg=TecAzul,fg=white}
\setbeamercolor{palette tertiary}{bg=TecAzulOscuro!80!black,fg=white}
\setbeamercolor{title}{fg=TecAzulOscuro,bg=TecFondoBloque}
\setbeamercolor{frametitle}{fg=TecAzulOscuro,bg=TecFondoBloque}

% Bloques personalizados elegantes
\setbeamercolor{block title}{bg=TecAzulOscuro,fg=white}
\setbeamercolor{block body}{bg=TecFondoBloque,fg=black}
\setbeamercolor{block title alerted}{bg=TecRojo,fg=white}
\setbeamercolor{block body alerted}{bg=TecRojo!8,fg=black}
\setbeamercolor{block title example}{bg=TecAzul,fg=white}
\setbeamercolor{block body example}{bg=TecAzul!8,fg=black}

% Redondear bordes y añadir sombra sutil
\setbeamertemplate{blocks}[rounded][shadow=true]
\setbeamertemplate{navigation symbols}{}

% Pie de página limpio con número de diapositiva
\setbeamertemplate{footline}{
  \leavevmode%
  \hbox{%
  \begin{beamercolorbox}[wd=.7\paperwidth,ht=2.25ex,dp=1ex,left]{author in head/foot}%
    \hspace*{2ex}\insertshortauthor~~(\insertshortinstitute) \hfill \insertshorttitle
  \end{beamercolorbox}%
  \begin{beamercolorbox}[wd=.3\paperwidth,ht=2.25ex,dp=1ex,right]{date in head/foot}%
    \insertshortdate{}\hspace*{2em}
    \insertframenumber{} / \inserttotalframenumber\hspace*{2ex}
  \end{beamercolorbox}}%
  \vskip0pt%
}

% Configuración de Listings de Python para visualización pedagógica de código
\lstset{
    language=Python,
    basicstyle=\ttfamily\footnotesize,
    keywordstyle=\color{TecAzulOscuro}\bfseries,
    stringstyle=\color{TecRojo},
    commentstyle=\color{TecGris}\itshape,
    showstringspaces=false,
    breaklines=true,
    frame=lines,
    rulecolor=\color{TecAzulOscuro!30},
    backgroundcolor=\color{TecFondoBloque},
    aboveskip=0.5em,
    belowskip=0.5em,
    literate={á}{{\'a}}1 {é}{{\'e}}1 {í}{{\'i}}1 {ó}{{\'o}}1 {ú}{{\'u}}1
             {Á}{{\'A}}1 {É}{{\'E}}1 {Í}{{\'I}}1 {Ó}{{\'O}}1 {Ú}{{\'U}}1
             {ñ}{{\~n}}1 {Ñ}{{\~N}}1 {¿}{{?`}}1 {¡}{{!`}}1
}

% Metadatos institucionales por defecto
\title[Modelación Estadística]{Modelación Estadística para Ciencia de Datos e Ingeniería}
\author[J. Castillo Colmenares]{Juliho David Castillo Colmenares}
\institute[Tec de Monterrey]{Tecnológico de Monterrey \\ Escuela de Ingeniería y Ciencias}
\date{\today}

% Comandos y macros estadísticas compartidas para las presentaciones Beamer (Metropolis)

% Conjuntos numéricos básicos
\newcommand{\R}{\mathbb{R}}
\newcommand{\N}{\mathbb{N}}
\newcommand{\Q}{\mathbb{Q}}
\newcommand{\Z}{\mathbb{Z}}
\newcommand{\set}[1]{\left\{ #1 \right\}}
\newcommand{\abs}[1]{\left|#1\right|}
\newcommand{\norm}[1]{\left\| #1 \right\|}

% Comandos estadísticos y probabilísticos del curso
\newcommand{\Var}{\operatorname{Var}}
\newcommand{\cov}{\operatorname{Cov}}
\newcommand{\corr}{\rho}
\newcommand{\comb}[2]{\begin{pmatrix} #1 \\ #2 \end{pmatrix}}
\newcommand{\s}{\sigma}
\newcommand{\particion}{\mathfrak{P}}
\newcommand{\card}[1]{\#\left( #1 \right)}
\newcommand{\seq}[3]{\set{#1}_{#2}^{#3}}
\newcommand{\E}{\mathbb{E}}
\newcommand{\Prob}{\mathbb{P}}

% Entornos pedagógicos y cajas en español académico natural
\newenvironment{discusion}{%
  \begin{alertblock}{Pregunta de Discusión}%
}{%
  \end{alertblock}%
}

\newenvironment{reflexion}{%
  \begin{alertblock}{Reflexión para el Aula}%
}{%
  \end{alertblock}%
}

\newenvironment{paradoja}{%
  \begin{alertblock}{Paradoja de Exploración}%
}{%
  \end{alertblock}%
}

\newenvironment{motivacion}{%
  \begin{exampleblock}{Motivación: Ciencia de Datos en la Práctica}%
}{%
  \end{exampleblock}%
}

\newenvironment{retoclase}{%
  \begin{block}{Reto Rápido en Equipo}%
}{%
  \end{block}%
}


\title[Estadísticos Z y t]{Sección 5.9: Estadísticos Z y t}
\subtitle{Estandarización con desviación conocida o desconocida}
\author[J. Castillo Colmenares]{Juliho Castillo Colmenares}
\institute{Tecnológico de Monterrey}
\date{}

\begin{document}
\begin{frame}[plain]\vspace{-0.8cm}\titlepage\end{frame}

\begin{frame}{Hoja de ruta --- capítulo 5}
  \footnotesize
  \begin{itemize}\itemsep=0.08cm
    \item \textbf{05.08} Conceptos fundamentales
    \item \textbf{05.09} Estadísticos Z y t \emph{(hoy)}
    \item \textbf{06--07} Intervalos y pruebas de hipótesis
  \end{itemize}
  \begin{block}{Objetivo didáctico}
    Construir $Z$ y $t$, distinguir cuándo se conoce $\sigma$ y leer cuantiles de
    la distribución $t$ de Student.
  \end{block}
\end{frame}

\begin{frame}{Configuración inferencial}
  \footnotesize
  La nota parte de un valor $A_0$ supuesto en la hipótesis nula y de una muestra
  aleatoria de tamaño $n$.
  \begin{block}{Cantidad observada}
    La media o medida muestral se denota $A$ y se compara con $A_0$.
  \end{block}
\end{frame}

\begin{frame}{Estadístico Z}
  \footnotesize
  \[
    Z=\frac{A-A_0}{\sigma/\sqrt n}.
  \]
  Aquí $\sigma$ es la desviación estándar poblacional conocida.
  \begin{alertblock}{Propósito}
    Convertir una variable normalizada en una normal estándar para consultar
    probabilidades y cuantiles tabulados.
  \end{alertblock}
\end{frame}

\begin{frame}{Caso Z: $\sigma$ conocida}
  \footnotesize
  Si la desviación estándar se conoce por experiencia previa, el estadístico Z
  usa directamente $\sigma/\sqrt n$.
  \begin{exampleblock}{Contexto de la nota}
    Los tiempos de entrega de una pizza ilustran una media cuya variabilidad
    poblacional ya está documentada.
  \end{exampleblock}
\end{frame}

\begin{frame}{Caso t: $\sigma$ desconocida}
  \footnotesize
  Cuando no existe un registro fiable de $\sigma$ o la muestra es pequeña,
  se sustituye por la desviación muestral $S$.
  \[
    t=\frac{A-A_0}{S/\sqrt n}.
  \]
  \begin{alertblock}{Diferencia}
    La incertidumbre adicional se refleja en la distribución $t$ de Student.
  \end{alertblock}
\end{frame}

\begin{frame}{Distribución de Student}
  \footnotesize
  La distribución $t$ fue descrita por William Sealy Gosset bajo el seudónimo
  Student. Es simétrica alrededor de cero y sus colas dependen de los grados de
  libertad.
  \begin{block}{Uso}
    Modela la estandarización cuando la desviación poblacional se estima con la
    propia muestra.
  \end{block}
\end{frame}

\begin{frame}{Grados de libertad y varianza}
  \footnotesize
  El parámetro `df` se denota $\nu$. Para $X\sim t_\nu$:
  \[
    \mu_X=0,\qquad \sigma_X^2=\frac{\nu}{\nu-2}
    \quad(\nu>2).
  \]
  Al crecer $\nu$, la distribución $t$ se aproxima a la normal estándar.
\end{frame}

\begin{frame}{Varianza muestral y estadístico t}
  \footnotesize
  La nota presenta
  \[
    S^2=\sum\frac{(A_i-A_0)^2}{n-1},
    \qquad t=\frac{A-A_0}{S/\sqrt n}.
  \]
  \begin{alertblock}{Lectura}
    El divisor $n-1$ y los grados de libertad conectan la estimación de $S$ con
    la forma de la distribución $t$.
  \end{alertblock}
\end{frame}

\begin{frame}{Puente computacional 1/4: elegir Z o t}
  \footnotesize
  \begin{tabular}{lll}
    \hline
    Información disponible & Estadístico & Distribución \\
    \hline
    $\sigma$ conocida & $Z=(A-A_0)/(\sigma/\sqrt n)$ & normal estándar \\
    $\sigma$ desconocida & $t=(A-A_0)/(S/\sqrt n)$ & $t_\nu$ \\
    \hline
  \end{tabular}
\end{frame}

\begin{frame}{Puente computacional 2/4: grados de libertad}
  \footnotesize
  Los grados de libertad determinan la cola de $t_\nu$. Con menos información
  sobre la desviación, las colas son más pesadas.
  \begin{block}{Límite}
    $t_\nu\to N(0,1)$ cuando $\nu$ crece.
  \end{block}
\end{frame}

\begin{frame}{Puente computacional 3/4: áreas}
  \footnotesize
  Para un valor observado $t$, la probabilidad a la derecha es
  \[
    P(T_\nu>t)=1-F_{t_\nu}(t).
  \]
  La simetría permite obtener también cuantiles izquierdos cambiando el signo.
\end{frame}

\begin{frame}{Puente computacional 4/4: cuantiles}
  \footnotesize
  La función percentil (ppf) satisface
  \[
    F_{t_\nu}(t_p)=p.
  \]
  Para un área derecha de $0.05$, se consulta $p=0.95$ y se usa la simetría
  para el valor negativo.
\end{frame}

\begin{frame}{Ejemplo resuelto 1/4 --- cuantiles con $\nu=9$}
  \footnotesize
  Consideremos una variable $t$ con $\nu=9$ grados de libertad. Se busca un
  valor con área derecha $0.05$ y área central sin sombrear $0.90$.
  \begin{block}{Planteamiento}
    Encontrar los dos límites simétricos.
  \end{block}
\end{frame}

\begin{frame}{Ejemplo resuelto 1/4 --- solución}
  \footnotesize
  El área izquierda del límite positivo es $0.95$:
  \[
    t_{0.95,9}\approx1.833,
    \qquad t_{0.05,9}\approx-1.833.
  \]
  \begin{alertblock}{Comprobación}
    $P(-1.833<T<1.833)=0.95-0.05=0.90$.
  \end{alertblock}
\end{frame}

\begin{frame}{Aplicación de lectura 2/4 --- prueba Z}
  \footnotesize
  Un proceso de entrega tiene desviación poblacional conocida. La media
  observada es $A$ y el valor de referencia es $A_0$.
  \begin{block}{Planteamiento}
    Estandarizar la diferencia con la fórmula Z de la sección.
  \end{block}
\end{frame}

\begin{frame}{Aplicación de lectura 2/4 --- interpretación Z}
  \footnotesize
  \[
    Z=\frac{A-A_0}{\sigma/\sqrt n}.
  \]
  Un valor positivo indica que $A$ está por encima de $A_0$; su magnitud se
  interpreta en una normal estándar.
\end{frame}

\begin{frame}{Aplicación de lectura 3/4 --- prueba t}
  \footnotesize
  Si $\sigma$ no está disponible, se calcula $S$ con la muestra y se usa
  \[
    t=\frac{A-A_0}{S/\sqrt n}.
  \]
  \begin{block}{Planteamiento}
    La incertidumbre sobre $\sigma$ debe reflejarse en $t_\nu$.
  \end{block}
\end{frame}

\begin{frame}{Aplicación de lectura 3/4 --- interpretación t}
  \footnotesize
  El valor observado se compara con cuantiles de $t_\nu$, no con cuantiles
  normales, especialmente cuando $n$ es pequeño.
  \begin{alertblock}{Resultado}
    Las colas más pesadas protegen contra la subestimación de la variabilidad.
  \end{alertblock}
\end{frame}

\begin{frame}{Aplicación de lectura 4/4 --- varianza muestral}
  \footnotesize
  La sección presenta $S^2$ con divisor $n-1$ y después forma el estadístico
  t usando $S$.
  \begin{block}{Planteamiento}
    Identificar qué cantidad cambia cuando aumenta la muestra.
  \end{block}
\end{frame}

\begin{frame}{Aplicación de lectura 4/4 --- conexión}
  \footnotesize
  Al crecer $n$, aumentan los grados de libertad, $S$ estima mejor a $\sigma$ y
  $t_\nu$ se acerca a $N(0,1)$.
  \begin{exampleblock}{Lectura integrada}
    La fórmula, la varianza muestral y la distribución de referencia forman un
    único procedimiento de estandarización.
  \end{exampleblock}
\end{frame}

\begin{frame}{Síntesis y transición}
  \footnotesize
  \begin{itemize}\itemsep=0.12cm
    \item Z usa $\sigma$ conocida y referencia normal.
    \item t usa $S$, grados de libertad y colas de Student.
    \item La simetría y la función percentil permiten leer áreas.
  \end{itemize}
  \begin{alertblock}{Siguiente paso}
    Estas herramientas se aplican en las secciones de intervalos y pruebas de
    hipótesis.
  \end{alertblock}
\end{frame}
\end{document}
